\chapter{Introducere}

Studiul relatiilor dintre active financiare reprezinta o problema importanta in analiza pietelor financiare. Evolutia unui sector economic poate influenta evolutia altor sectoare prin mecanisme macroeconomice sau prin transmiterea informatiei.

In aceasta lucrare sunt utilizate metode statistice moderne pentru a studia structura dependentei dintre sectoare ale pietei bursiere.

Abordarea utilizata combina modele vector autoregresive cu metode de regularizare pentru estimarea matricii de precizie. Rezultatul este reprezentat sub forma unei retele financiare care descrie dependentele conditionale dintre variabile.

Analiza de retea permite identificarea sectoarelor centrale si a structurii generale a pietei.